\chapter{Narrative Infrastructure}
\label{ch:narrative-infrastructure}

\epigraph{Who controls the past controls the future. Who controls the present controls the past.}{George Orwell, \textit{Nineteen Eighty-Four}}

\noindent
Beliefs do not float freely.
They travel through infrastructure: schools, media organizations, platforms, think tanks, entertainment industries, professional training systems, and religious institutions.
Just as an economy requires physical infrastructure (roads, ports, power grids) to move goods, a doxonomic system requires \emph{narrative infrastructure} to move beliefs from production to population.

This chapter models narrative infrastructure as a directed network of transformation functions with measurable capacity, throughput, concentration, and vulnerability.
The infrastructure is not neutral: its topology determines which narratives can travel far and which are stillborn.


\section{Infrastructure as Transformation Network}
\label{sec:ch7-transformation}

\begin{definition}[Narrative Infrastructure]
\label{def:narrative-infrastructure}
\index{narrative infrastructure}
A \emph{narrative infrastructure} $\narr$ is a directed weighted graph $G = (V, E, w)$ where:
\begin{itemize}[nosep]
  \item each vertex $v \in V$ is an institution with processing function $\iota_v: \R_{\geq 0} \times \narspace \to \narspace$,
  \item each directed edge $(u, v) \in E$ carries a flow weight $w_{uv} \in [0,1]$ representing the fraction of $u$'s output that serves as input to $v$,
  \item the network has source nodes (funders injecting resources) and sink nodes (audience segments absorbing narratives).
\end{itemize}
\end{definition}

The key insight is that narrative infrastructure is not a flat landscape.
It has topology: some paths are wide and well-maintained (the pipeline from elite think tanks through prestige media to policy debate), while others are narrow or blocked (the path from grassroots organizations to national attention).
The infrastructure determines \emph{which} narratives are structurally advantaged, independent of their truth or popularity.

\begin{example}[The Conservative Narrative Pipeline]
\label{ex:conservative-pipeline}
\index{narrative pipeline}
The American conservative movement built one of the most effective narrative pipelines in modern history, documented in \textcite{mayer2016} and \textcite{mirowski2009}:
\[
  \text{Donors} \xrightarrow{0.8} \text{Foundations} \xrightarrow{0.9} \text{Think tanks} \xrightarrow{0.5} \text{Op-ed pages} \xrightarrow{0.6} \text{Cable news} \xrightarrow{0.7} \text{Public}
\]
The flow weights indicate the estimated fraction of output that propagates to the next stage.
The end-to-end propagation probability is $\prod_i w_i = 0.8 \times 0.9 \times 0.5 \times 0.6 \times 0.7 \approx 0.15$, seemingly low, but applied to hundreds of reports per year across dozens of institutions, it produces a constant stream of narrative output reaching the public through high-credibility channels.

A parallel pipeline runs through academia:
\[
  \text{Donors} \xrightarrow{0.7} \text{Endowed programs} \xrightarrow{0.6} \text{Textbooks} \xrightarrow{0.8} \text{Curricula} \xrightarrow{0.95} \text{Students}
\]
This pipeline has lower throughput per year but vastly higher persistence ($\persistence{} \sim$ decades).
\end{example}

\begin{remark}[Infrastructure vs.\ Content]
\label{rem:infra-content}
A crucial distinction: narrative infrastructure is about the \emph{channels}, not the \emph{messages}.
The same infrastructure that carries ``markets are efficient'' could, in principle, carry ``markets are prone to failure.''
But infrastructure is not neutral: institutions built to produce one type of narrative develop staff, expertise, donor relationships, and organizational culture that make producing alternative narratives costly.
A think tank funded for 30 years to argue against regulation does not easily pivot to arguing for it.
Infrastructure has inertia.
\end{remark}


\section{Institutional Types and Their Infrastructure Roles}
\label{sec:ch7-types}

Different institutions occupy structurally different positions in the narrative infrastructure.
Their roles can be classified by where they sit in the production chain and what transformation they perform.

\begin{definition}[Infrastructure Role Taxonomy]
\label{def:role-taxonomy}
\index{infrastructure roles}
\begin{description}[nosep]
  \item[Generators] produce original narrative content: research universities, think tanks, investigative journalists, research labs.
  \item[Amplifiers] increase the reach of existing narratives without substantially transforming them: broadcast media, social media platforms, wire services, syndication networks.
  \item[Translators] transform narratives from one register to another: textbook authors (research $\to$ pedagogy), popularizers (academic $\to$ public), lobbyists (policy research $\to$ legislative language).
  \item[Validators] attach credibility to narratives: peer review systems, professional associations, credentialing bodies, editorial boards.
  \item[Gatekeepers] control access to high-value channels: editors, algorithm designers, curriculum committees, conference organizers.
\end{description}
\end{definition}

A single institution may perform multiple roles.
The \textit{New York Times} is simultaneously a generator (original reporting), an amplifier (wire service redistribution), a translator (making complex events accessible), a validator (the prestige of publication implies quality), and a gatekeeper (deciding what is ``fit to print'').
This multi-role character gives major media institutions outsized structural power.

\begin{center}
\begin{tikzpicture}[
  every node/.style={font=\small},
  box/.style={draw, thick, rounded corners, minimum width=3cm, minimum height=0.8cm, align=center},
]
\matrix[row sep=0.5cm, column sep=0.3cm] {
  \node[box, fill=doxoblue!8] {Education}; &
  \node[font=\small, text width=7.5cm] {\textbf{Generator + Translator + Validator.} High persistence ($\sim$decades), high credibility, low throughput per unit time. Shapes foundational priors. Most consequential long-run infrastructure.}; \\
  \node[box, fill=doxoblue!8] {Prestige Media}; &
  \node[font=\small, text width=7.5cm] {\textbf{Generator + Amplifier + Gatekeeper.} Medium persistence ($\sim$months), high credibility, medium throughput. Sets the daily agenda and frames for elites and opinion leaders.}; \\
  \node[box, fill=doxoblue!8] {Social Media}; &
  \node[font=\small, text width=7.5cm] {\textbf{Amplifier + (sometimes) Generator.} Very low persistence ($\sim$hours), low credibility, very high throughput. Amplifies, fragments, and accelerates. Algorithms act as invisible gatekeepers.}; \\
  \node[box, fill=doxoblue!8] {Think Tanks}; &
  \node[font=\small, text width=7.5cm] {\textbf{Generator + Validator.} High credibility in policy circles, low direct public reach. Produces the ``studies show'' citations that legitimate policy narratives.}; \\
  \node[box, fill=doxoblue!8] {Entertainment}; &
  \node[font=\small, text width=7.5cm] {\textbf{Translator + Amplifier.} High reach, low perceived persuasive intent. Normalizes through narrative representation, not argument. Operates at layers 3--5.}; \\
  \node[box, fill=doxoblue!8] {Religious Institutions}; &
  \node[font=\small, text width=7.5cm] {\textbf{Generator + Validator + Community.} Very high persistence, high credibility within congregations, medium reach. Unique access to moral framing and identity formation.}; \\
  \node[box, fill=doxoblue!8] {Professional Training}; &
  \node[font=\small, text width=7.5cm] {\textbf{Translator + Validator.} Converts academic knowledge into operational assumptions. MBA programs, law schools, medical training embed anthropological regimes into professional practice.}; \\
};
\end{tikzpicture}
\end{center}


\section{Throughput and Capacity}
\label{sec:ch7-throughput}

\begin{definition}[Narrative Throughput]
\label{def:throughput}
\index{narrative throughput}
The \emph{throughput} of an institution $v$ is the quantity
\[
  \Theta(v) = \funding{v} \cdot \reach{v} \cdot \persistence{v}
\]
measuring the total belief-exposure-hours produced per unit time.
The \emph{capacity} $\Theta^{\max}(v)$ is the maximum throughput achievable given institutional constraints (staffing, infrastructure, brand).
\end{definition}

Throughput combines three dimensions: how much is spent, how many people are reached, and how long the effect lasts.
This composite measure allows comparison across radically different channel types.

\begin{example}[Throughput Comparison]
\label{ex:throughput-comparison}
Consider three channels, each receiving \$1M/year:
\begin{center}
\begin{tabular}{lcccc}
\toprule
\textbf{Channel} & $\reach{}$ & $\persistence{}$ (years) & $\Theta$ (M\$·person·yr) \\
\midrule
University economics course & 0.001 & 30 & 0.03 \\
Cable news segment & 0.05 & 0.02 & 0.001 \\
Textbook chapter & 0.005 & 20 & 0.1 \\
\bottomrule
\end{tabular}
\end{center}
The textbook has 100 times the throughput of cable news, despite reaching far fewer people per day, because its persistence is three orders of magnitude higher.
This is why doxonomic strategists who understand the mathematics invest in education and publishing, not just media.
\end{example}

\begin{remark}[Why Throughput Is Not Impact]
\label{rem:throughput-impact}
Throughput measures exposure, not belief change.
A high-throughput channel with low credibility (social media) may produce less actual belief change than a low-throughput channel with high credibility (a mentor's advice).
The full impact depends on the internalization operator $\mathcal{I}$ (Chapter~\ref{ch:belief-production-chain}), which weights exposure by credibility, resistance, and audience characteristics.
Throughput is a necessary but not sufficient condition for doxonomic effect.
\end{remark}


\section{Network Topology and Structural Advantage}
\label{sec:ch7-topology}

The topology of narrative infrastructure creates structural advantages and disadvantages for different narratives.

\begin{definition}[Narrative Pathway]
\label{def:pathway}
\index{narrative pathway}
A \emph{narrative pathway} from source $s$ to audience segment $t$ is a directed path $s = v_0, v_1, \ldots, v_l = t$ in the infrastructure graph $G$.
The \emph{pathway efficiency} is
\[
  \eta(s \to t) = \prod_{i=0}^{l-1} w_{v_i, v_{i+1}} \cdot \credibility{v_i}
\]
measuring the fraction of the original narrative that reaches the audience, weighted by the credibility accumulated along the way.
\end{definition}

\begin{proposition}[Structural Advantage]
\label{prop:structural-advantage}
\index{structural advantage}
A narrative $n_1$ has a \emph{structural advantage} over narrative $n_2$ if the maximum pathway efficiency from any source to the target audience is higher for $n_1$:
\[
  \max_{s} \eta_{n_1}(s \to t) > \max_{s} \eta_{n_2}(s \to t)
\]
Structural advantage arises from the infrastructure, not from the content.
A narrative with better-connected institutional sponsors, more pathways to high-credibility amplifiers, and shorter routes to the target audience has a structural advantage regardless of its truth value.
\end{proposition}

\begin{example}[Structural Advantage: Deregulation vs.\ Public Ownership]
\label{ex:structural-advantage}
The narrative ``deregulation promotes efficiency'' has extensive infrastructure: well-funded think tanks (Heritage, Cato, AEI) connected to prestigious op-ed pages (\textit{WSJ}, \textit{FT}), cable news panels, congressional testimony channels, and economics curricula.
The maximum pathway efficiency to U.S.\ policymakers is high.

The narrative ``public ownership can be efficient'' has minimal infrastructure: a few academic advocates, limited think-tank support, almost no connections to major op-ed pages or cable news.
The maximum pathway efficiency to the same policymakers is low.

This structural difference exists \emph{independently} of the empirical merits of either claim.
It is a property of the infrastructure, not of the evidence.
\end{example}


\section{Concentration Metrics}
\label{sec:ch7-concentration}

\begin{definition}[Epistemic Herfindahl Index]
\label{def:ehhi-infra}
\index{epistemic Herfindahl index}
The \emph{epistemic Herfindahl index} for narrative infrastructure is
\[
  \ehhi = \sum_{v \in V} \left(\frac{\Theta(v)}{\sum_{u \in V} \Theta(u)}\right)^2
\]
Values near $1$ indicate monopoly in narrative production; values near $1/|V|$ indicate deconcentration.
\end{definition}

\begin{conjecture}[Concentration and Capture]
\label{conj:concentration-capture}
If $\ehhi > \bar{h}$ for a critical threshold $\bar{h}$, then a single institution or small coalition has the \emph{capacity} to achieve common-sense capture (Definition~\ref{def:common-sense}) on propositions within its narrative domain, provided sufficient funding and absent strong counter-institutions, resistant audiences, or contradicting lived experience.
The threshold $\bar{h}$ is not a universal constant; it depends on audience receptivity, competing credibility structures, and proposition-content fit.
\end{conjecture}

\begin{example}[Concentration in U.S.\ Media]
\label{ex:us-concentration}
As of 2020, five companies (Comcast/NBCUniversal, Disney, ViacomCBS, WarnerMedia, News Corp/Fox) controlled approximately 90\% of U.S.\ media consumption \autocite{bagdikian2004}.
If we weight by throughput (reach $\times$ persistence), the $\ehhi$ for U.S.\ national media is approximately 0.20--0.25, near the threshold typically associated with high concentration in market economics.

But this understates doxonomic concentration, because the five companies share substantial overlap in:
\begin{itemize}[nosep]
  \item advertising dependence (same advertisers, same incentive structure),
  \item sourcing patterns (same wire services, same government press conferences),
  \item professional culture (journalism schools produce a shared worldview),
  \item ownership class interests (all publicly traded, all dependent on deregulatory environment).
\end{itemize}
Effective doxonomic concentration is higher than $\ehhi$ alone suggests, because the ``competing'' outlets share structural filters (Chapter~\ref{ch:prehistory}, Section~\ref{sec:ch3-herman-chomsky}).
\end{example}

\begin{definition}[Effective Concentration]
\label{def:effective-concentration}
\index{effective concentration}
The \emph{effective concentration} adjusts $\ehhi$ for narrative correlation:
\[
  \ehhi_{\text{eff}} = \ehhi + (1 - \ehhi) \cdot \bar{\rho}
\]
where $\bar{\rho} \in [0,1]$ is the average pairwise narrative correlation between outlets (measured by text similarity, framing overlap, or source sharing).
When all outlets produce identical narratives ($\bar{\rho} = 1$), $\ehhi_{\text{eff}} = 1$ regardless of market structure: a perfectly correlated oligopoly is doxonomically equivalent to a monopoly.
\end{definition}


\section{Infrastructure Resilience and Vulnerability}
\label{sec:ch7-resilience}

\begin{definition}[Infrastructure Vulnerability]
\label{def:vulnerability}
\index{infrastructure vulnerability}
A narrative infrastructure is \emph{vulnerable} at node $v$ if removing $v$ significantly reduces the maximum pathway efficiency for a class of narratives.
The \emph{vulnerability index} of $v$ is
\[
  \nu(v) = 1 - \frac{\max_s \eta_{-v}(s \to t)}{\max_s \eta(s \to t)}
\]
where $\eta_{-v}$ is the pathway efficiency with $v$ removed.
High $\nu(v)$ means $v$ is a critical bottleneck.
\end{definition}

\begin{example}[Bottleneck Institutions]
\label{ex:bottleneck}
In many countries, a small number of editorial gatekeepers control access to the prestige media pipeline.
If three editors at major newspapers decide that a topic is not newsworthy, the pathway from think-tank research to public attention is effectively blocked.
These editors have high vulnerability indices: removing them from the network (or changing their editorial stance) would significantly alter which narratives reach the public.

Similarly, a curriculum committee that approves textbooks for a state with millions of students is a high-vulnerability node.
Texas's State Board of Education, which approves textbooks used by 5.5 million students, has historically been a bottleneck node with outsized doxonomic influence on national textbook content (publishers produce one version for Texas and sell it nationally).
\end{example}

\begin{proposition}[Infrastructure Fragility Under Concentration]
\label{prop:fragility-concentration}
Highly concentrated infrastructure ($\ehhi$ high) is more efficient for narrative transmission but more fragile: disruption of a single high-throughput node has a larger system-wide effect.
Formally, if $\Theta(v) / \sum_u \Theta(u) > 1/2$ for some $v$, then $\nu(v) > 1/2$ and the infrastructure is critically dependent on $v$.
\end{proposition}

This creates a design tension: concentrated infrastructure is powerful but fragile; distributed infrastructure is resilient but less efficient.
Chapter~\ref{ch:epistemic-democracy} considers what institutional design best balances efficiency and resilience.


\section{Infrastructure Construction and Path Dependence}
\label{sec:ch7-construction}

Narrative infrastructure does not appear overnight.
Building it requires sustained investment over years or decades.

\begin{example}[Building Conservative Infrastructure]
\label{ex:building-infrastructure}
The conservative narrative infrastructure in the United States was consciously constructed over 40 years \autocite{mayer2016, phillipsfein2009}:
\begin{itemize}[nosep]
  \item \textbf{1947}: Mont P\`{e}lerin Society founded (intellectual seed).
  \item \textbf{1971}: Lewis Powell memo urges American business to invest in ideological infrastructure.
  \item \textbf{1973}: Heritage Foundation founded (\$250K initial; \$80M+ budget by 2020).
  \item \textbf{1977}: Cato Institute founded.
  \item \textbf{1980s}: Olin, Scaife, and Bradley foundations dramatically increase grants to conservative law-and-economics programs at universities.
  \item \textbf{1990s}: Fox News launched (1996), completing the broadcast amplification layer.
  \item \textbf{2000s}: Digital expansion (Daily Wire, Turning Point USA, PragerU).
\end{itemize}
Total estimated investment: several billion dollars over four decades.
The result: a vertically integrated pipeline from academic seed-planting through policy production, media amplification, and political mobilization.

No comparable left-of-center infrastructure was built during the same period \autocite{mudge2018}.
This infrastructure asymmetry, not a difference in the quality of ideas, explains much of the rightward shift in American policy common sense between 1980 and 2020.
\end{example}

\begin{principle}[Infrastructure Path Dependence]
\label{prin:path-dependence}
\index{path dependence!infrastructure}
Once built, narrative infrastructure is self-reinforcing: institutions produce alumni who staff other institutions, generate citations that boost credibility, create demand for their own output, and develop donor relationships that ensure continued funding.
Replicating an established infrastructure from scratch requires investment of comparable scale, which is why infrastructure asymmetries persist across decades.
\end{principle}


\section{Platform Infrastructure: The Algorithmic Layer}
\label{sec:ch7-platforms}

Digital platforms add a new layer to narrative infrastructure: the algorithmic layer, which sits between institutional production and audience reception.

\begin{model}[Algorithmic Infrastructure]
\label{mod:algorithmic-infra}
\index{algorithmic infrastructure}
A platform algorithm is a function $A: \narspace \times \mathcal{U} \to [0,1]$ mapping each narrative-user pair $(n, j)$ to a display probability $A(n, j)$.
The algorithm determines the effective exposure matrix:
\[
  e_{jk}^{\text{eff}} = e_{jk} \cdot A(n_k, j)
\]
The platform's objective function (typically engagement maximization) determines $A$ endogenously:
\[
  A^* = \arg\max_A \sum_{j,k} A(n_k, j) \cdot \text{engagement}(n_k, j)
\]
Because high-engagement content tends toward moral outrage, identity threat, and simplistic narratives, engagement maximization systematically distorts the effective narrative infrastructure \autocite{zuboff2019}.
\end{model}

\begin{remark}[Platforms as Invisible Gatekeepers]
\label{rem:invisible-gatekeepers}
Traditional gatekeepers (editors, curriculum committees) are visible and accountable.
Algorithmic gatekeepers are invisible and optimized for a private objective (engagement/revenue).
This means the most consequential infrastructure decisions (which narratives reach which audiences) are made by systems that are neither transparent nor democratically accountable.
From a doxonomic perspective, platform algorithms may be the single most powerful gatekeeper institution in the contemporary information environment, yet they are the least understood and least regulated.
\end{remark}


\section{Notes and References}
\label{sec:ch7-notes}

Media concentration is studied in \textcite{bagdikian2004} and \textcite{mcchesney1999}.
Comparative media systems analysis is developed in \textcite{hallin2004}.
The role of think tanks in narrative production is documented in \textcite{mayer2016}, \textcite{mirowski2009}, and \textcite{phillipsfein2009}.
Platform infrastructure is analyzed in \textcite{zuboff2019} and \textcite{benkler2018}.
The network model of media draws on \textcite{jackson2008}.
The Herfindahl index originates in \textcite{herfindahl1950}.
For public-interest alternatives, see \textcite{benkler2006}.
On the propaganda model's structural filters as infrastructure properties, see \textcite{herman1988}.
The concept of narrative correlation adjusting effective concentration is new to this text but draws on the literature on media parallelism in \textcite{hallin2004}.
On the construction of conservative intellectual infrastructure, see \textcite{mayer2016} and \textcite{phillipsfein2009}; on the absence of a comparable left infrastructure, see \textcite{mudge2018}.


\section{Exercises}

\begin{exercise}
Compute $\ehhi$ for U.S.\ national television news using publicly available audience share data.
Then estimate the average narrative correlation $\bar{\rho}$ (using a qualitative assessment of framing similarity) and compute $\ehhi_{\text{eff}}$.
How much does the correction change the concentration assessment?
\end{exercise}

\begin{exercise}
Model the narrative pipeline from the Atlas Network to public opinion on regulation in a specific country.
Draw the infrastructure graph, estimate the flow weights $w_{uv}$ at each stage, and compute the end-to-end pathway efficiency.
Where is the pipeline most likely to fail?
\end{exercise}

\begin{exercise}
Prove that $\ehhi \geq 1/|V|$ with equality if and only if all institutions have equal throughput.
\end{exercise}

\begin{exercise}
A new social media platform enters the narrative infrastructure with $\reach{} = 0.4$, $\credibility{} = 0.15$, $\persistence{} = 0.05$.
How does this change the overall $\ehhi$?
Does it increase or decrease the conditions for common-sense capture?
(Consider both the direct effect on concentration and the indirect effect on narrative correlation.)
\end{exercise}

\begin{exercise}
\label{ex:ch7-bottleneck}
Identify the three highest-vulnerability nodes in the narrative infrastructure of your country (using the vulnerability index $\nu(v)$).
For each, describe what would happen if that node were captured by a single funder or if it ceased to function.
\end{exercise}

\begin{exercise}
Compare the conservative narrative pipeline (Example~\ref{ex:building-infrastructure}) with the infrastructure available to labor movements or environmental movements.
Draw both pipeline graphs.
Where are the structural asymmetries?
What would it cost to close the infrastructure gap?
\end{exercise}

\begin{exercise}
\label{ex:ch7-algorithm}
Model a platform algorithm that maximizes engagement.
Show that if outrage-inducing content generates higher engagement, the algorithm will systematically amplify narratives with high emotional valence and low nuance.
What does this imply for the narrative entropy $H$ (Definition~\ref{def:narrative-entropy}) of the platform's output?
\end{exercise}

\begin{exercise}
The throughput comparison (Example~\ref{ex:throughput-comparison}) suggests that textbooks are 100 times more cost-effective than cable news.
Why, then, do doxonomic funders invest heavily in cable news?
Give at least three reasons why the throughput calculation does not capture the full strategic picture.
\end{exercise}
